% Source PDF: source/普通高考/2015/2015全国2文(贵州,甘肃,青海,西藏,黑龙江,吉林,海南,宁夏,内蒙古,新疆,云南,广西,辽宁).pdf, page 1.
% PDF embedded-bitmap check: no image objects; the frequency histogram is vector content.
% Source bars (score groups 40--50 through 90--100) have frequency densities
% 0.010, 0.020, 0.030, 0.020, 0.015, 0.005 respectively.
% solid segments: six histogram bars and coordinate axes; dashed segments: guides
% at 0.005, 0.010, 0.015, 0.020 and 0.030. Labels avoid the bars and axes.

\begin{tikzpicture}
\begin{axis}[
    width=8.4cm,
    height=4.8cm,
    xmin=35,
    xmax=105,
    ymin=0,
    ymax=0.046,
    axis lines=none,
    xtick=\empty,
    ytick=\empty,
    clip=false,
    enlargelimits=false,
    scale only axis,
]
    % Frequency-density guide lines read from the original histogram.
    \draw[dashed,line width=0.35pt] (axis cs:35,0.005) -- (axis cs:100,0.005);
    \draw[dashed,line width=0.35pt] (axis cs:35,0.010) -- (axis cs:100,0.010);
    \draw[dashed,line width=0.35pt] (axis cs:35,0.015) -- (axis cs:100,0.015);
    \draw[dashed,line width=0.35pt] (axis cs:35,0.020) -- (axis cs:100,0.020);
    \draw[dashed,line width=0.35pt] (axis cs:35,0.030) -- (axis cs:100,0.030);

    % Six equal-width score groups and their frequency densities.
    \addplot[/tikz/ybar interval,mark=none,draw=black,fill=white,line width=0.45pt] coordinates {
      (40,0.010)
      (50,0.020)
      (60,0.030)
      (70,0.020)
      (80,0.015)
      (90,0.005)
      (100,0)
    };

    % Axes and labels.
    \draw[->,black,line width=0.45pt] (axis cs:35,0) -- (axis cs:104,0);
    \draw[->,black,line width=0.45pt] (axis cs:35,0) -- (axis cs:35,0.042);
    \node[anchor=south,inner sep=0pt] at (axis cs:67.5,0.046) {A 地区用户满意度评分的频率分布直方图};
    \node[anchor=north east,inner sep=0pt] at (axis cs:34.3,-0.0004) {$O$};
    \node[anchor=south west,inner sep=0pt] at (axis cs:35.7,0.037)
        {\(\dfrac{\text{频率}}{\text{组距}}\)};
    \node[anchor=west,inner sep=0pt] at (axis cs:104.2,0.0002) {满意度评分};
    \node[anchor=east,inner sep=1pt] at (axis cs:34.5,0.005) {0.005};
    \node[anchor=east,inner sep=1pt] at (axis cs:34.5,0.010) {0.010};
    \node[anchor=east,inner sep=1pt] at (axis cs:34.5,0.015) {0.015};
    \node[anchor=east,inner sep=1pt] at (axis cs:34.5,0.020) {0.020};
    \node[anchor=east,inner sep=1pt] at (axis cs:34.5,0.025) {0.025};
    \node[anchor=east,inner sep=1pt] at (axis cs:34.5,0.030) {0.030};
    \node[anchor=east,inner sep=1pt] at (axis cs:34.5,0.035) {0.035};
    \node[anchor=east,inner sep=1pt] at (axis cs:34.5,0.040) {0.040};
    \node[anchor=north,inner sep=0pt,font=\normalsize] at (axis cs:40,-0.001) {40};
    \node[anchor=north,inner sep=0pt,font=\normalsize] at (axis cs:50,-0.001) {50};
    \node[anchor=north,inner sep=0pt,font=\normalsize] at (axis cs:60,-0.001) {60};
    \node[anchor=north,inner sep=0pt,font=\normalsize] at (axis cs:70,-0.001) {70};
    \node[anchor=north,inner sep=0pt,font=\normalsize] at (axis cs:80,-0.001) {80};
    \node[anchor=north,inner sep=0pt,font=\normalsize] at (axis cs:90,-0.001) {90};
    \node[anchor=north,inner sep=0pt,font=\normalsize] at (axis cs:100,-0.001) {100};
\end{axis}
\end{tikzpicture}
